%=========== ΛΥΣΗ - 22ο ΘΕΜΑ Δ ===========
\begin{thema}{Δ}
\begin{erwthma}
\item Η $f$ είναι παραγωγίσιμη στο $\mathbb{R}$ και για κάθε $x\in\mathbb{R}$ είναι
\begin{align*}
f'(x)
&=2+\frac{2x}{x^2+1}\\
&=\frac{2(x^2+x+1)}{x^2+1}.
\end{align*}
Επειδή
\[
x^2+1>0
\]
και
\[
x^2+x+1=\left(x+\frac{1}{2}\right)^2+\frac{3}{4}>0,
\]
έχουμε
\[
f'(x)>0
\]
για κάθε $x\in\mathbb{R}$. Άρα η $f$ είναι γνησίως αύξουσα και επομένως συνάρτηση 1-1.

Επίσης
\[
\lim_{x\to-\infty}f(x)=-\infty
\quad\text{και}\quad
\lim_{x\to+\infty}f(x)=+\infty.
\]
Άρα
\[
f(\mathbb{R})=\mathbb{R},
\]
οπότε η $f$ αντιστρέφεται και το πεδίο ορισμού της αντίστροφης είναι
\[
{D_{f^{-1}}=\mathbb{R}}.
\]

Επειδή η $f$ είναι γνησίως αύξουσα, έχουμε
\begin{align*}
f(x^2-x)>f(x+3)
&\Leftrightarrow x^2-x>x+3\\
&\Leftrightarrow x^2-2x-3>0\\
&\Leftrightarrow (x+1)(x-3)>0.
\end{align*}
Άρα το σύνολο λύσεων είναι
\[
{x\in(-\infty,-1)\cup(3,+\infty)}.
\]

\item Θεωρούμε τη συνάρτηση
\[
h(x)=f(x)-\frac{1}{x},
\quad x\in(0,1].
\]
Η $h$ είναι παραγωγίσιμη και
\[
h'(x)=f'(x)+\frac{1}{x^2}>0.
\]
Άρα η $h$ είναι γνησίως αύξουσα στο $(0,1]$. Επιπλέον,
\[
\lim_{x\to0^+}h(x)
=\lim_{x\to0^+}\left(f(x)-\frac{1}{x}\right)
=-\infty
\]
και
\[
h(1)=f(1)-1=1+\ln2>0.
\]
Επομένως, από το Θεώρημα Ενδιάμεσων Τιμών, υπάρχει $x_0\in(0,1)$ τέτοιο ώστε $h(x_0)=0$. Η ρίζα αυτή είναι μοναδική, διότι η $h$ είναι γνησίως αύξουσα.

Άρα η εξίσωση
\[
f(x)=\frac{1}{x}
\]
έχει ακριβώς μία ρίζα στο $(0,1)$.

\item Έχουμε
\[
f(0)=0
\quad\text{και}\quad
f'(0)=2.
\]
Επομένως η εφαπτομένη της $C_f$ στο σημείο $O(0,0)$ έχει εξίσωση
\[
y-f(0)=f'(0)(x-0)
\Leftrightarrow
{y=2x}.
\]

Για κάθε $x\in\mathbb{R}$ είναι
\begin{align*}
f''(x)
&=\left(\frac{2x}{x^2+1}\right)'\\
&=\frac{2(x^2+1)-4x^2}{(x^2+1)^2}\\
&=\frac{2(1-x^2)}{(x^2+1)^2}.
\end{align*}
Στο διάστημα $[0,1]$ έχουμε
\[
f''(x)\geq0,
\]
με $f''(x)>0$ για $x\in[0,1)$. Άρα η $f$ είναι αυστηρά κυρτή στο $[0,1]$.

Επομένως η $C_f$ βρίσκεται πάνω από την εφαπτομένη της στο $O$. Για κάθε $x\in(0,1]$ ισχύει
\[
f(x)>2x.
\]
Ολοκληρώνοντας κατά μέλη στο $[0,1]$, παίρνουμε
\[
\int_0^1f(x)\d x
>\int_0^1 2x\d x
=1.
\]
Άρα
\[
{\int_0^1f(x)\d x>1}.
\]

\item Είναι
\begin{align*}
g(x)
&=xf(x)\\
&=2x^2+x\ln(x^2+1).
\end{align*}
Στο διάστημα $[0,1]$ έχουμε $x\geq0$ και $f(x)\geq f(0)=0$, άρα $g(x)\geq0$. Επομένως το ζητούμενο εμβαδόν είναι
\begin{align*}
E
&=\int_0^1g(x)\d x\\
&=\int_0^1\left[2x^2+x\ln(x^2+1)\right]\d x\\
&=\frac{2}{3}
+\frac{1}{2}\int_1^2\ln u\,\d u\\
&=\frac{2}{3}
+\frac{1}{2}\left[u\ln u-u\right]_1^2\\
&=\frac{2}{3}+\ln2-\frac{1}{2}\\
&={\ln2+\frac{1}{6}}.
\end{align*}
\end{erwthma}
\end{thema}
%=========== ΤΕΛΟΣ ΛΥΣΗΣ - 22ο ΘΕΜΑ Δ ===========
